\section{SCENARIO \#1: "KERNSTOWN": March 22-26, 1862
}

\begin{scenariodata}
\textbf{GAME LENGTH:} 5 Turns\\
(WEATHER: March Column)\\
\textbf{FIRST PHASE:} Confederates
\end{scenariodata}

\begin{scenariodata}
\textbf{UNION:}\\
\textbf{RESOURCE FACTOR:} 8\\
\textbf{INITIAL PLACEMENT:}\\
Hex B-37 (Franklin) : (3I], 3I, 1A, 1C, (1F, 1S)\\
Hex B-25 (Moorefield) : (3I, 1A], (1F, 1S)\\
Hex A-16 (Romney) : [5I, 1A, 1C], (1F, 1S)\\
Hex N-20 (Winchester) : 19I, 5A, 2C, \textbf{and adjacent :} (Shields, 1F,
4S, 2W)\\
B\&O : [18I, 2A, 3C], (1F, 2S)\\
Hex W-21 : 18I, 3A, 2C, \textbf{and adjacent :} (Banks, 1W, 1S)\\
\textbf{WITHDRAWAL HEXES:}\\
N-7, S-10, Z-13, CC-21 through CC-33\\
\textbf{SUPPLY/WAGON ENTRY HEXES:}\\
N-7, S-10, Z-13, A-16, A-11, A-13, CC-21 through CC-33\\
No additional Fortification counters.
\end{scenariodata}

\begin{scenariodata}
\textbf{CONFEDERATES:}\\
\textbf{RESOURCE FACTOR:} 5\\
\textbf{INITIAL PLACEMENT:}\\
Hex M-29 (Woodstock) : 10I, 5A, \textbf{and adjacent :} (Jackson, 2S, 3W)\\
Within 6 hexes of Hex M-29 (Woodstock) : 3C, 1H, (Ashby)\\
Hex E-46 (McDowell) : \textbf{and adjacent :} 7I, 1A, (1S)\\
Hex O-51 (Staunton) : 1H, (1F, 1S)\\
Hex BB-53 (Charlottesville) : 1I, (1F)\\
(Partisans: McNeill, Gilmor)\\
\textbf{WITHDRAWAL HEXES:}\\
CC-47 through CC-58\\
\textbf{SUPPLY/WAGON ENTRY HEXES:}\\
CC-47 through CC-58, O-51, M-50\\
No additional Fortification counters.
\end{scenariodata}

\begin{scenariodata}
\textbf{SPECIAL \#1:} If Winchester is held by the Confederates (any type of
Combat unit) at the end of the game, it is worth 10 Victory Points.\\
\textbf{SPECIAL \#2:} Confederate intelligence was less than perfect during
this early campaign. Allow the Union Player to place his units in their
Stacking Boxes upside-down. The Confederate Player cannot examine these units
unless adjacent.\\
\textbf{SPECIAL \#3:} Jackson had asked for additional troops, and the Union
command was unsure if he had gotten them. To simulate this uncertainty, the
Confederate Player may introduce an additional 7I, 1A to the game. These are
placed with the rest of the main army near Woodstock, and the Union Player is
not informed if these are present or not. If used, these raise the Confederate
Resource Factor by one. Adjust Victory Point totals accordingly at the end of
the game if these troops have been used.
\end{scenariodata}

\begin{scenariodata}
\textbf{VICTORY CONDITIONS:}\\
\textbf{ADVANCED GAME:} Union wins if they have 5 or more lead in total
Victory Points. Confederates win if they have 15 or more lead in total Victory
Points. All other results are a draw.\\
\textbf{BASIC GAME:} Same as the Advanced Game, only "spot" the Confederates
10 Victory Points at the start of the game.\\
\textbf{HISTORICAL RESULTS:} Draw.
\end{scenariodata}

\begin{scenariodata}
\textbf{PLAYTESTER'S COMMENTS:} The Union Player must be careful not to keep
more units on the Mapboard than needed, a chancy guessing game. A rigorous
recon effort is needed by the cavalry to develop the actual Rebel strength and
positions. Pitch into the Confederates for attrition whenever possible. The
Confederates have a good option for attacks on Moorefield and Romney, but don't
leave the Pike unguarded. Maneuvers towards Harper's Ferry may also panic the
Union, but basically, try to avoid action, except for possibly a last-minute
rush on Winchester or Harper's Ferry.
\end{scenariodata}

Although repulsed at Kernstown, Jackson's advance seemed menacing enough to the
Union authorities to persuade them to order still more units into the
Shenandoah region. Jackson managed to elude these forces for a month while
awaiting fresh Confederate reinforcements. By the end of April, 1862, the
division of Richard S. "Baldy" Ewell (1817-1872) was approaching the Blue Ridge
gaps. Ewell was an ex-dragoon officer with an excellent record in the Mexican
War, and as an Indian-fighter. At first doubting his new commander's sanity,
Ewell was eventually to give Jackson his whole-hearted support. A driving
commander, Ewell's concept of war is best summed up in his statement, "The road
to glory cannot be followed with much baggage."

Ewell's arrival coincided with the opening moves of the Union spring offensive
from West Virginia. These Union troops were commanded by the famous
"Pathfinder," John C. Fremont (1813-1890), another political appointee who had
been the first Republican nominee for President in 1856. McDowell had been
captured, Staunton was threatened, and still another Federal division was
passing through Winchester on the way to join Fremont's army...
